%% kapa: exam=1 round=24 year=2013 period=2 form=A subject=accounting subject_ko=회계학 seq=38 number=38
%% source: Q-Net 감정평가사 기출문제 (cst00309) artlSeq=5073802 file=extracted/1008/5073802/2013년 제24회 감정평가사 1차시험 기출문제/2교시(A형).hwp
%% conversion: kordoc 4.24.51 -> tools/gen1.py (auto); review=proofread-sonnet; flags=none
\begin{question}{38}
(주)감평은 이익중심점인 A사업부와 B사업부를 운영하고 있다. A사업부가 생산하는 열연강판의 변동제조원가와 고정제조원가는 각각 톤당 ￦2,000과 톤당 ￦200이며, 외부 판매가격과 판매비는 각각 톤당 ￦3,000과 톤당 ￦100이다. 현재 B사업부가 열연강판을 외부에서 톤당 ￦2,600에 구입하여 사용하고 있는데, 이를 A사업부로부터 대체받을 것을 고려하고 있다. A사업부는 B사업부가 필요로 하는 열연강판 수요를 충족시킬 수 있는 유휴생산능력을 보유하고 있으며, 사내대체하는 경우 판매비가 발생하지 않을 것이다. A사업부가 사내대체를 수락할 수 있는 최소사내대체가격은?
\begin{choices}
  \item ￦2,000
  \item ￦2,100
  \item ￦2,200
  \item ￦2,600
  \item ￦3,000
\end{choices}
\end{question}
\answer{1}
